\chapter{What comes after attention}
\label{ch:post-attention}

\newthought{Sparse attention} of any flavour is still attention. The
selection function is new; the dot-product-and-softmax routing inside the
selected set is unchanged. Several research programmes in 2026 are asking
the more aggressive question: \emph{is the entire attention pattern
the right computational primitive at all?}

This is speculative. None of the alternatives below has a frontier-quality
deployment yet. The point of including them is to remind the reader that
the post-transformer landscape might extend well past where this book has
mapped it.

\section{Diffusion language models}

Discrete diffusion models \citep{lou2024discrete} apply the
denoising-diffusion machinery (which has dominated image generation since
2022) to text. Instead of predicting the next token autoregressively,
they iteratively refine a draft of the entire output, all positions in
parallel.

\paragraph{Why this could matter.} The autoregressive loop is the
dominant cost of LLM inference: generating 1000 tokens takes 1000
sequential forward passes. A diffusion LM that produces a fluent draft
in 20 refinement steps is potentially 50$\times$ faster at the same
quality.

\paragraph{Why it hasn't won yet.} Quality. Discrete diffusion language
models in 2026 are competitive with mid-size autoregressive transformers
but not with frontier models. The gap may close with scale; nobody is
sure.

\section{JEPA and the world-model camp}

Yann LeCun's Joint Embedding Predictive Architectures
\citep{lecun2022jepa} argue that LLMs train on the wrong objective: they
predict surface symbols (tokens) instead of latent abstractions. A JEPA
predicts the embedding of the future, not the future itself, and uses
that embedding-prediction task as the basis for representation learning.

\paragraph{The bet.} If you can train a world model that predicts
\emph{abstract} state rather than concrete tokens, you should get a model
that plans and reasons better. Whether you can train such a world model
purely on text — without the multimodal grounding LeCun has emphasised —
is contested.

\paragraph{Status.} Early. JEPA is a research direction, not a deployed
product, as of mid-2026.

\section{Programmatic / neuro-symbolic systems}

A different bet: instead of making models bigger, make them \emph{call
out} to deterministic computation (search engines, code interpreters,
databases) and treat the LLM as a \emph{controller} of those tools rather
than the source of all answers. Agents in this style are already common;
the open question is how much of the LLM's parametric knowledge can
eventually be moved to external tools, and whether that lets the LLM
itself be much smaller.

\section{Energy-based and continuous-state models}

A small but persistent line of work pursues sequence models that maintain
a continuous, low-dimensional state and update it via energy
minimisation rather than autoregression. The connection to control
theory and physics-informed models is strong; the connection to
production NLP is, so far, weak.

\section{Honest closing}

A book like this is by definition out of date the moment it ships. The
$\bigO(n^2)$ wall is real, the engineering response is mature, and the
architectural response is still in flux. SSA is one credible bet on the
hardest quadrant of that response. Whether it is \emph{the} bet is going
to be settled by independent reproduction, by training-cost economics,
and by what happens when the next generation of long-context benchmarks
makes everyone start over again.

The interesting part of the field is the part where the answer is not yet
obvious. We are squarely in that part right now.

\keyidea{Attention has been the centre of language modelling for a
decade. There is no law that says it has to stay there. The honest
position — for engineers, not marketers — is to bet on whichever
architecture serves the customer cheapest at the quality the customer
needs, and to keep reading the papers.}

\fillme{FILL-20-01}{Sidebar: diffusion language models for text people}{%
  ${\sim}400$-word sidebar walking a transformer-trained engineer through
  discrete diffusion: forward process (corrupt tokens with masking or
  uniform noise), reverse process (learn to denoise), the trade-off vs
  autoregression (parallel decoding but harder training objective). Cite
  Lou et al.\ 2024 and Hoogeboom et al.\ 2021 (multinomial diffusion).
  End with one sentence on Inception Labs / Mercury demos as the 2025
  proof-of-concept that this works at deployable speed.}{%
  Lou et al.\ 2024; Hoogeboom et al.\ 2021; Inception Labs Mercury blog.}{%
  ${\sim}400$ words; one schematic equation OK; no figure required.}

\fillme{FILL-20-02}{Sidebar: what JEPA actually predicts}{%
  ${\sim}400$-word sidebar with a clear word-picture of the JEPA training
  objective: encoder \(f\) maps an input to a latent embedding; predictor
  \(g\) takes a context view's latent and predicts a target view's latent;
  loss is in latent space, not pixel / token space. Contrast with
  autoregressive next-token loss. End with one sentence on V-JEPA
  (video-JEPA) results and one on the open question of whether JEPA-style
  objectives transfer to text.}{%
  LeCun 2022 position paper; Assran et al.\ 2023 I-JEPA; Bardes et al.\
  2024 V-JEPA.}{%
  ${\sim}400$ words; one schematic equation; no figure required.}
